% Unnumbered diagnostic. Period one and a nonminimal anchor; no counting claim.
\subsection*{Order generation on a curve does not imply the small incidence bound}

Let $k$ be algebraically closed, let $d\ge1$, and work on
$X=\mathbb P^1_k$.  Put
\[
 G=\mathcal O_X(-1)\oplus\mathcal O_X(1)^{d-1},\qquad
 H=\mathcal O_X^d,\qquad
 B=\mathcal End(G),\qquad E=\mathcal Hom(G,H).
\]
Use precomposition for the right $B$-action.  The order $B$ is a split
Azumaya algebra of ordinary rank $d^2$, and $E$ is locally a free
right $B$-module of rank one.  Nevertheless the exact minimum number
of global right-$B$ generators of $E$ is
\[
                                  d+1.                 \tag{B.PG.1}
\]
Thus actual global order generation does not give the $s+1$ incidence
bound when $s=1$; for $d\ge2$ the minimum exceeds two.

Indeed
\[
 E\simeq\mathcal O_X(1)^d\oplus
                         \mathcal O_X(-1)^{d(d-1)}.
\]
Every global map $G\to H$ therefore vanishes on the summand
$\mathcal O(1)^{d-1}$.  It is specified by one column of $d$ linear
forms, or equivalently by a map $\mathcal O(-1)\to\mathcal O^d$.
For $b$ such maps $\varphi_1,\ldots,\varphi_b$, their right-$B$ span
equals $E$ precisely when their image columns span $H$ at every point.
In local frames this is the elementary statement that matrices with
given image columns generate the regular matrix module under right
multiplication if and only if those columns span $k^d$.  Hence global
generation by the $b$ maps is equivalent to a surjection
\[
                   \mathcal O_X^b\twoheadrightarrow
                                      \mathcal O_X(1)^d. \tag{B.PG.2}
\]
For $b<d$ this is impossible by fibre rank.  For $b=d$ its determinant
would be a nowhere vanishing section of $\mathcal O_X(d)$, which is
impossible because a nonzero homogeneous form of positive degree on
$\mathbb P^1$ has a zero over $k$.

For $b=d+1$, use homogeneous coordinates $[x:y]$ and the $d$ by $d+1$
matrix having $x$ in entry $(i,i)$, $y$ in entry $(i,i+1)$, and zero
elsewhere.  If $x\ne0$, its first $d$ columns have determinant $x^d$.
If $x=0$, then $y\ne0$ and its last $d$ columns have determinant $y^d$.
It therefore defines the required surjection (B.PG.2) in every
characteristic.  This proves (B.PG.1), as well as actual global
$B$-generation of $E$.

The example also realizes the elementary fibre diagnostic: in
$E(P)=\operatorname{Mat}_{d,d}(k)$ the image of the global section
space consists of first-column matrices.  Its right-$B(P)$ span is
the whole module, but $b$ scalar choices can span at most $b$ image
directions at a point.  Global geometry strengthens the fibre lower
bound $b\ge d$ to the exact $d+1$.

This is a period-one, nonminimal-anchor example: a rank-one anchor
would suffice for the split Brauer class.  It asserts no canonical
surface-order or fixed-unit-pole realization, and no counterexample
to optimized rank-period anchors or to a hereditary representation
count.  Its exact purpose is to prevent replacing the evaluation
hypothesis of the small-generator incidence argument by order
generation alone.  All calculations above are internal; no new
external source or promoted statement is claimed.
